\section{MI-19: a rank-two counterexample to the subset inequality for the \texorpdfstring{$q$}{q}-permanent}
\label{sec:mi19}
For $A=(a_{ij})\in\M_n$ and $q\in\R$, use the inversion convention
\[
 \operatorname{per}_q(A)=\sum_{\sigma\in S_n}q^{\ell(\sigma)}\prod_{i=1}^n a_{i,\sigma(i)},\qquad
 \ell(\sigma)=\#\{(i,j):i<j,\ \sigma(i)>\sigma(j)\}.
\]
For a subset $S\subseteq\{1,\ldots,n\}$, define the restricted sum
\[
 R_{q,S}(A)=\sum_{\substack{\sigma\in S_n\\\sigma(S)=S}}
 q^{\ell(\sigma)}\prod_{i=1}^n a_{i,\sigma(i)}.
\]
The target is $\operatorname{per}_q(A)\geq R_{q,S}(A)$ for Hermitian positive semidefinite $A$ and $0\leq q\leq1$. This is the subset formulation recorded in \cite[Conjecture 3, equation (5.1)]{Fonseca2018}.

\begin{theorem}
The subset inequality is false for a real positive semidefinite matrix of order four and rank two. A counterexample is
\[
 A=\begin{pmatrix}
 170&1&167&170\\
 1&1&-2&1\\
 167&-2&173&167\\
 170&1&167&170
 \end{pmatrix},\qquad q=\frac78,\qquad S=\{2\}.
\]
\end{theorem}
\begin{proof}
Positive semidefiniteness and rank two follow from
\[
 A=X^TX,\qquad X=\begin{pmatrix}13&0&13&13\\1&1&-2&1\end{pmatrix}.
\]
Expansion over the $24$ permutations gives
\begin{align*}
 \operatorname{per}_q(A)
 &=115600q^6+4886140q^5+9568969q^4+4712758q^3\\
 &\hspace{1.5em}-199231q^2+4886140q+4999700.
\end{align*}
The six permutations fixing the second index give
\[
 R_{q,\{2\}}(A)=4999700q^5+9482260q^4+4741130q^3+4741130q+4999700.
\]
Consequently the difference is the small polynomial
\begin{align*}
 \operatorname{per}_q(A)-R_{q,\{2\}}(A)
 &=q(q+1)\bigl(115600q^4-229160q^3+315869q^2\\
 &\hspace{8em}-344241q+145010\bigr).
\end{align*}
At $q=7/8$ this equals
\[
 -\frac{3235575}{16384}<0.
\]
All calculations are integer polynomial arithmetic. The accompanying verifier independently enumerates the permutations with the stated inversion convention.
\end{proof}

\begin{remark}[Positive definite variants and scope]
The strict negative gap persists for $A+\varepsilon I_4$ for every sufficiently small positive $\varepsilon$, by continuity of the two polynomial expressions. Thus the obstruction is not dependent on allowing a singular input. No counterexample to the separate monotonicity problem MI-18 is asserted.
\end{remark}
